\appendix

\section{Systems with algebraic Hamiltonians}\label{App0}

First, let us remark that Theorem~\ref{thm:MoRa} also holds for a
general Poisson system \cite{akn}. When the Hamiltonian function is algebraic but not meromorphic,
we cannot apply this theorem directly. One solution is to find an extension of the phase space (by including additional variables) in~such a way, that the original Hamiltonian lifts to a meromorphic one, and the extended system is Hamiltonian
with respect to a degenerate Poisson bracket, which reproduces
the original problem.

The construction below is a modification of that given in
\cite{alg}, where the reader will find more details and proofs. Let us
consider an $n$ degrees of freedom Hamiltonian system with canonical
coordinates $q,p\in\mathbb{C}^n$, with algebraic Hamiltonian $H(q,p)$
such that $H(q,p)=K(q,p,u)$, where~$u$ is algebraic over
$\mathbb{C}(q)$ with minimal polynomial $P(u)\in\mathbb{C}(q)[u]$, and
$K(q,p,u)\in \mathbb{C}(q,p,u)$ is a~rational function of its
arguments $x=(q,p,u)\in\mathbb{C}^{2n+1}$. We introduce the following
system
\begin{equation}
 \label{eq:dK}
 \dot x = J(x) \nabla_xK(x),
\end{equation}
where $J(x)$ is $(2n+1)\times(2n+1)$ matrix of the form
\begin{equation*}
 J(x):=\begin{bmatrix}
 0 & \mathbbm{1}_n & 0 \\
 -\mathbbm{1}_n & 0 & \frac{1}{\partial_u P}\nabla_q P \\
 0 & -\frac{1}{\partial_u P} \nabla_q P& 0
 \end{bmatrix}\!,
\end{equation*}
with $\mathbbm{1}_n$ equal to the $n\times n$ identity matrix. It defines the Poisson bracket
\begin{equation*}
 \{f,g\}(x) := (\nabla_xf(x))^{\rm T} J(x)\nabla_xg(x),
\end{equation*}
where $f$ and $g$ are smooth functions. The rank of matrix $J(x)$ is
$2n$ and the only Casimir function of the bracket is $P(u)$.
\begin{Lemma}
{\sloppy If $(q(t),p(t),u(t))$ is a solution of equations~\eqref{eq:dK} with
 $P(u(t))=0$, then $(q(t),p(t))$ is a solution of Hamilton's
 equations
 \begin{equation*}
 %\label{eq:has}
 \dot q= \nabla_q H(q,p), \qquad
 \dot p=-\nabla_p H(q,p).
 \end{equation*}}
\end{Lemma}
We omit the proof, as it is rather direct, and ask the interested reader to follow the explanation and steps given in \cite[Section~2]{alg}. This lemma gives us what is needed, that is we reproduce the original system as a~Hamiltonian one with respect to a degenerate Poisson structure defined
by rational matrix $J(x)$, and with rational Hamiltonian function~$K(x)$.

The above general considerations justify the meromorphic assumptions
of the Morales--Ramis theory, but of course for practical purposes the
calculations can be performed in the original coordinates.
